\section{Introduction}

The classical global-regularity problem for the three-dimensional incompressible Navier--Stokes equations sits inside a long PDE line running from Leray and Hopf weak solutions through the monographs of Ladyzhenskaya and Temam, with partial regularity and continuation criteria due to Caffarelli--Kohn--Nirenberg, Prodi, Serrin, Beale--Kato--Majda, Escauriaza--Seregin--Sver\'ak, and Koch--Tataru isolating adjacent regularity mechanisms \cite{Leray1934,Hopf1951,Ladyzhenskaya1969,Temam2001,CaffarelliKohnNirenberg1982,Prodi1959,Serrin1962,BealeKatoMajda1984,EscauriazaSereginSverak2003,KochTataru2001}.

We consider the classical three-dimensional incompressible Navier-Stokes system
\begin{equation}
\partial_t u + (u\cdot\nabla)u + \nabla p = \nu\Delta u,\qquad \nabla\cdot u = 0.
\end{equation}
The domain is \([0,T)\times\mathbb{R}^3\), and the initial data are smooth, compactly supported, and divergence-free. The theorem proved here stays on that classical Euclidean surface throughout. In particular, results obtained only for regularized equations are treated as diagnostic context unless they are recovered on the classical system without theorem-critical modification.

The paper is organized around four bridge propositions. The first is a scale barrier that prevents indefinite transfer into dangerous high frequencies. The second is the monotone functional \(Q(t)\), which records the global energy/enstrophy layer on the same classical surface and keeps the original framework honest about what compactness is allowed to consume. The third is a compactness upgrade strong enough to pass the nonlinear term on the classical equation. The fourth is the curvature-based regularity bridge, now stated in the export package as a pure Euclidean gradient-control theorem on the declared classical surface. The present manuscript reconstructs that source-grounded route explicitly and keeps the Euclidean bridge theorem separate from the remaining conditional route theorem rather than conflating them.

The organizational choice is therefore simple: keep the original four-way loop explicit, state the exact conditional route theorem that the present bridge package supports, and keep title, abstract, theorem, proof, and conclusion aligned with that honest claim boundary. The realized classical carrier and exact pullback are recorded as same-surface support infrastructure, while Bridge IV itself is stated and proved separately in Euclidean form on the declared surface.
